\chapter{Notation and Planning Templates}
\label{app:notation}

\section{Notation Table}

\begin{center}
\begin{tabularx}{\textwidth}{@{}l X@{}}
\toprule
Symbol & Meaning \\
\midrule
$s$ & behavioural state \\
$\FI(s)$ & framework inertia of state $s$ \\
$\Reward(s)$ & immediate reward rate of the current state \\
$\Switch(s)$ & cognitive switching cost of leaving the current state \\
$\Env(s)$ & environmental reinforcement of the current state \\
$\DV(s_0 \to s_1)$ & activation barrier for a transition \\
$\Will(t)$ & volitional energy available at time $t$ \\
$\Leverage(t)$ & effective leverage of free will at time $t$ \\
$P$ & prior preparation that lowers future barrier terms \\
\bottomrule
\end{tabularx}
\end{center}

\section{Planning Templates}

\begin{templatebox}
\textbf{Session-design card}
\begin{enumerate}[nosep]
  \item Target outcome for the week:
  \item Next session that serves the outcome:
  \item Exact first action:
  \item Objects that must already be placed:
  \item Window in which the session should begin:
  \item Rule for ending the session:
\end{enumerate}
\end{templatebox}

\begin{templatebox}
\textbf{Failure review card}
\begin{enumerate}[nosep]
  \item Which state won?
  \item Which barrier term was underestimated?
  \item Which cue restored the old state?
  \item Which smaller entry action should replace the current one?
  \item What will be changed before the next window arrives?
\end{enumerate}
\end{templatebox}

\section{Appendix Use}

The appendix exists for reuse. If later versions add more formal proofs,
empirical citations, or domain-specific case banks, those modules can still
share the same notation and planning templates. That keeps the manuscript
coherent as it grows.

